\begin{studiumworked}{M05. Respuesta local de un conducto subsónico · Intermedio}
\textbf{Enunciado.} Una sección académica tiene Mach 0.40. Su área disminuye un 1\,\%. Estime el cambio relativo de rapidez y densidad mediante el modelo diferencial.

\textbf{Datos.} Mach: $0.4$; Variación área: $-0.01$.

\textbf{Incógnitas.} Magnitudes solicitadas y cierre físico.

\textbf{Hipótesis.} Modelo académico con hipótesis del capítulo. Datos originales; no representan hardware ni ensayo.

\textbf{Modelo físico.} Conservación y relaciones constitutivas declaradas.
\[dA/A=(M^2-1)du/u\]
\[du/u=(dA/A)/(M^2-1),\quad d\rho/\rho=-M^2du/u\]
\textbf{Desarrollo.}
\begin{enumerate}
\item Despejar la velocidad da $(-0.01)/(0.4^2-1)=0.01190476$. Es un aumento de 1.190476\,\%.
\item La densidad cambia por $-0.4^2(0.01190476)=-0.001904762$, es decir, disminuye aproximadamente 0.1904762\,\%.
\item Sumar las tres variaciones logarítmicas da cero: la menor área se compensa con rapidez mayor y densidad algo menor.
\end{enumerate}
\textbf{Resultados.} $du/u\simeq0.01190476$; $d\rho/\rho\simeq-0.001904762$.

\textbf{Interpretación.} El resultado corresponde a las hipótesis indicadas; su interpretación requiere conservar fronteras y unidades.

\textbf{Comprobación.} La continuidad linealizada cierra exactamente; para una variación finita el error es de segundo orden.

\textbf{Límites.} Sólo una estimación local, lejos de Mach uno; no determina régimen de todo el conducto.
\end{studiumworked}
